\section{Conclusion, Limitations, and Future Directions}
\label{conclusion}
% In this work, we presented \textit{EnCoMP}, an enhanced covert maneuver planning framework that combines multi-modal perception and offline reinforcement learning to enable autonomous robots to navigate safely and efficiently in complex environments. Our approach introduces several key contributions, including a multi-modal perception pipeline that generates high-fidelity cover, threat, and height maps, an offline reinforcement learning algorithm that learns robust navigation policies from real-world datasets, and an effective integration of perception and learning components for informed decision-making. However, there are a few limitations to our current approach. First, the performance of \textit{EnCoMP} relies on the quality and diversity of the real-world dataset used for training. While we have collected a substantial amount of data, there may still be scenarios or environments that are not well-represented in the dataset. Second, our approach assumes that the environment is static during navigation, which may not always hold in real-world scenarios. Dynamic obstacles or changing terrain conditions could impact the effectiveness of the learned policy. These limitations indicate the need for further research and development to enhance the robustness and adaptability of \textit{EnCoMP} in complex, real-world environments.

% Future research directions include extending \textit{EnCoMP} to handle dynamic environments by incorporating online adaptation mechanisms into the reinforcement learning framework. Additionally, we plan to investigate the integration of active perception techniques, such as active mapping or exploration, to allow the robot to actively gather information about the environment during navigation. An exciting opportunity for future research is to scale up EnCoMP to groups or teams of robots, enabling collaborative and coordinated navigation in complex environments. This direction opens up possibilities for multi-robot systems to tackle challenging missions more efficiently and effectively. Furthermore, we plan to validate the performance of \textit{EnCoMP} on a wider range of real-world terrains and scenarios, including more diverse vegetation types, weather conditions, and mission objectives.


In this work, we presented \textit{EnCoMP}, an enhanced covert maneuver planning framework that combines LiDAR-based perception and offline reinforcement learning to enable autonomous robots to navigate safely and efficiently in complex environments. Our approach introduces several key contributions, including an advanced perception pipeline that generates high-fidelity cover, threat, and height maps, an offline reinforcement learning algorithm that learns robust navigation policies from real-world datasets, and an effective integration of perception and learning components for informed decision-making. Additionally, we introduced the Adaptive Threat-Aware Visibility Estimation (ATAVE) algorithm, which dynamically assesses and mitigates potential threats in real-time, enhancing the robot's situational awareness and decision-making capabilities.
However, there are a few limitations to our current approach. First, the performance of \textit{EnCoMP} relies on the quality and diversity of the real-world dataset used for training. While we have collected a substantial amount of data, there may still be scenarios or environments that are not well-represented in the dataset. Second, our approach assumes that the environment is static during navigation, which may not always hold in real-world scenarios. Dynamic obstacles or changing terrain conditions could impact the effectiveness of the learned policy. These limitations underscore the need for further research and development to enhance the robustness and adaptability of \textit{EnCoMP} in complex, real-world environments.

Future research directions include extending \textit{EnCoMP} to handle dynamic environments by incorporating online adaptation mechanisms into the reinforcement learning framework. Additionally, we plan to investigate the integration of active perception techniques, such as active mapping or exploration, to allow the robot to actively gather information about the environment during navigation. Another opportunity for future research is to scale up \textit{EnCoMP} to groups or teams of robots, enabling collaborative and coordinated navigation in complex environments. This direction opens up possibilities for multi-robot systems to tackle challenging missions more efficiently and effectively. Furthermore, we plan to validate the performance of \textit{EnCoMP} on a wider range of real-world terrains and scenarios, including more diverse vegetation types, weather conditions, and mission objectives.

\section*{Acknowledgment}
This work has been partially supported by U.S. Army Grant \texttt{\#W911NF2120076} and ONR Grant \texttt{\#N00014-23-1-2119}